\subsection{特征标的正交关系}

我们保留第 5、6 小节的记号。回顾 \(Z(K[G])\) 由中心函数组成。

我们用公式定义一个从 \(K[G] \times K[G]\) 到 K 的对称双线性映射：

\[
\langle f, f ^ {\prime} \rangle_ {\mathrm{G}} = | \mathrm{G} | ^ {- 1} \sum_ {g \in \mathrm{G}} f _ {g} f _ {g ^ {- 1}} ^ {\prime}\tag{28}
\]

对每个 \(f = \sum f_{g}g\) 与 \(f' = \sum f_{g}'g\)（都属于 K{[}G{]}）成立。我们有 \(\langle f, f' \rangle_{\mathrm{G}} = |\mathrm{G}|^{-2}\tau (ff')\)。

命题 4（特征标的正交关系）

对 \(\lambda\) 与 \(\mu\)（在 \(\hat{\mathrm{G}}\) 中），我们有 \(\langle \chi_{\lambda},\chi_{\mu}\rangle_{\mathrm{G}} = \delta_{\lambda \mu}\)。

这是关系式 (20) 与 (23) 当自同态 u 与 v 取为恒同时的特殊情形。

推论. —— 设 \(\pi\) 与 \(\pi'\) 是 G 的有限维线性表示。在域 K 中，我们有

\[
\langle \chi_ {\pi}, \chi_ {\pi^ {\prime}} \rangle_ {\mathrm{G}} = (\dim_ {\mathrm{K}} \operatorname{Hom} _ {\mathrm{G}} (\pi , \pi^ {\prime})) \cdot 1.\tag{29}
\]

我们首先设 \(\pi\) 与 \(\pi'\) 是单表示。向量空间 \(\mathrm{Hom}_{\mathrm{G}}(\pi, \pi')\) 的维数依 \(\pi\) 与 \(\pi'\) 是否同构而为 1 或 0（舒尔引理，VIII, 第 47 页，命题 2）。在这一情形，公式 (29) 由命题 4 得出。

在一般情形，表示 \(\pi\)（相应地，\(\pi'\)）是单表示 \(\pi_{1},\ldots,\pi_{m}\)（相应地，\(\pi_{1}',\ldots,\pi_{n}'\)）的直和。空间 \(\operatorname{Hom}_{\mathrm{G}}(\pi,\pi')\) 同构于诸空间 \(\operatorname{Hom}_{\mathrm{G}}(\pi_{i},\pi_{j}')\)（\(1\leqslant i\leqslant m\)，\(1\leqslant j\leqslant n\)）的直和，并且我们有

\[
\chi_ {\pi} = \chi_ {\pi_ {1}} + \dots + \chi_ {\pi_ {m}}, \quad \chi_ {\pi^ {\prime}} = \chi_ {\pi_ {1} ^ {\prime}} + \dots + \chi_ {\pi_ {n} ^ {\prime}}.
\]

由线性性，公式 \((29)\) 的证明可化为单表示的情形。

